\section{Related Work}
\label{sec:related_work}
% Discuss existing approaches to failure detection and membership management in distributed systems. Compare
% heartbeat-based detection, timeout-based failure detectors, gossip-style membership protocols, and
% centralized monitoring. Explain which ideas influenced your design.
\subsection{Failure detection}

\subsubsection{\texorpdfstring{$\phi$}{phi} Accrual Calculation}

Hayashibara et al.~\cite{hayashibaraAccrualFailureDetector2004} proposed the \texorpdfstring{$\phi$}{phi} accrual
failure detector to improve the flexibility of failure detection. Instead of returning a boolean result, such as
``alive'' or ``unreachable'', the detector outputs a continuous suspicion level. This allows the system or developer to
interpret the monitoring information using different thresholds depending on the required trade-off between fast
detection and false-positive reduction. A lower threshold can detect failures more quickly, while a higher threshold
can reduce incorrect failure decisions. Experimental results from the original work show that the $\phi$ accrual
failure detector performs comparably to other adaptive failure detection mechanisms such as failure detector of Chen et
al \cite{chenQualityServiceFailure}. and that of Bertier et al. \cite{bertierImplementationPerformanceEvaluation2002a},
while providing greater flexibility for applications \cite{hayashibaraAccrualFailureDetector2004}.

\subsubsection{SWIM protocol}

Traditional heartbeating protocols may also suffer from scalability problems and false-positive failure
detection\cite{guptaScalableEfficientDistributed2001}. If every node monitors every other node, the number of heartbeat
messages increases rapidly as the system grows. In addition, if a node misses a heartbeat only because of temporary
network problems, it may incorrectly mark another node as failed. Das et al.~\cite{dasSWIMScalableWeaklyconsistent2002}
proposed SWIM to address these issues by separating failure detection from membership dissemination and using efficient
peer-to-peer probing. In each protocol period, a node selects one target node from its membership and sends a direct
ping. If the target responds, it is considered alive. If not, helper nodes perform indirect probes; a response through
any helper confirms that the target is still reachable.

A node is considered failed only when both the direct probe and the indirect probes fail. SWIM also disseminates
membership changes using an infection-style, or gossip-style, approach, allowing updates about joins, leaves,
suspicions, and failures to spread through the system with low per-node message overhead.

\subsection{System Design}
System architecture is another important design consideration. A common approach for coordinating distributed
applications is to introduce a dedicated coordination service, such as ZooKeeper
\cite{huntZooKeeperWaitfreeCoordination}. ZooKeeper is a replicated coordination service consisting of multiple
servers, one of which acts as the leader. Clients may connect to any server in the ensemble, where read operations are
served locally while write operations are forwarded to the leader, ordered, and replicated across the remaining
servers.

This architecture provides strong consistency guarantees, high availability, and good scalability for read-dominated
workloads. However, it also introduces a logically centralized coordination layer that must be maintained separately
from the application nodes. This increases deployment and operational complexity and makes membership management
dependent on a dedicated coordination ensemble.

For this reason, our system adopts a decentralized design instead. Rather than relying on an external coordination
service, membership management, failure detection, and topology maintenance are distributed across the participating
nodes themselves.

\subsection{Contributions}
\label{subsec:contribution}

The main contributions are as follows:

\begin{itemize}

  \item A novel neighbor discovery protocol that orchestrates node integration while preserving the structural properties of a
        k-regular graph overlay: FINDING NEMO.

  \item A hybrid version of SWIM and \texorpdfstring{$\phi$}{phi} accrual failure detector: hybrid SWIMPHI.

  \item A new self-healing topology maintainance technique that automatically restores the overlay toward a k-regular graph
        state following node failures or topology disruptions: CHUALANH.

  \item Implementation of the framework including core components above in Java.

  \item Direct evaluation of FINDING NEMO, CHUALANH, and the overall system performance via a dashboard that enables the
        observation of node behavior and the collection of experimental metrics for evaluation.

\end{itemize}
